\section{Design Principles}\label{sec:3}

The prototype described in \secref{sec:4} and the interface described in \secref{sec:5} instantiate four design principles that we derive here from prior work in human--{AI} interaction, explainability, and trust calibration.
The principles are not claims about the system's behavior; they are commitments against which the system's design is judged.
\secref{sec:7} returns to each principle and reports the empirical property that operationalizes it.

\leadpara{Principle G1: Transparency of state}
A user must be able to see, before any recommendation is computed, what the system has parsed from their document and what it will compare.
The principle is grounded in the observation that explanations of opaque model output are difficult to evaluate when the user cannot see the inputs to the model~\cite{amershi2019guidelines,liao2020questioning}.
In our system, the candidate portal exposes the parsed fields of a resume as a review form before any snapshot is registered; the employer portal exposes the parsed requirements of a job description as a similar review form.
A user who disagrees with a parsed field can edit it; the system does not rank on the basis of half-edited forms.
This is the transparency property the system commits to: the data on which a ranking is computed is the same data the user last confirmed.

\leadpara{Principle G2: Faithful, component-level explanations}
When the system returns a ranked list, every list item should carry an explanation that names the input fields that contributed to its position.
The principle follows the {XAI} literature that ties explanation quality to the question a user is actually asking~\cite{miller2019explanation,ehsan2022operationalizing} and to recent work that warns against explanations that are internally consistent but decoupled from the components that produced the score~\cite{joshi2021real}.
Our explanations are bound to the same composite channels that produced the score; we evaluate their faithfulness against the contributions of individual channels in \secref{sec:7}.

\leadpara{Principle G3: User control at consequential points}
Three actions in the system are consequential for the user: applying to a job, shortlisting a candidate, and contacting a candidate.
None of them happens implicitly.
Each requires a user gesture; the system records the action and updates local state but does not act on the user's behalf.
This commitment is the user-control counterpart of the trust-calibration literature~\cite{lee2004trust,parasuraman1997humans,shneiderman2020human} and of the mixed-initiative tradition that assigns the system a supporting role~\cite{horvitz1999principles,amershi2014power,kulesza2015principles}.
A useful interaction is one in which the user can always inspect what the system did, undo it if it was wrong, and decide what to do next.

\leadpara{Principle G4: Multi-perspective decision support}
A ranking is rarely a single number in the user's mind; it is a composite of independent considerations---text similarity, skill overlap, experience tier, compensation fit, and remote policy in our case.
The system exposes these as separable components of the score rather than collapsing them into a single value, so that a user who weights remote policy more heavily than experience tier can read both components in the explanation panel.
The principle is grounded in the recommendation literature that warns against over-reliance on a single signal~\cite{ziegler2005improving,tintarev2015explaining} and in the interactive-machine-learning tradition that exposes intermediate signals to the user for inspection and revision~\cite{kulesza2015principles,amershi2014power}.

\leadpara{How the principles relate}
The four principles are not independent.
G1 supports G2: a user who can see what the system has parsed can evaluate an explanation that names those fields.
G2 supports G3: a user who sees a component-level reason for a ranking is in a position to act on the components they disagree with.
G3 supports G4: when users retain control over the final action, exposing the components of a score is informative rather than threatening.
We refer back to this dependency in \secref{sec:8} when we discuss the design space our system does not address.
